% !TeX spellcheck = en_US
\documentclass[11pt]{article}
\usepackage{exerciseHead}

%%% ENCODING, FONT & LOCALIZATION
\usepackage[utf8]{inputenc} % use unicode input-encoding (allows unicode characters in the .tex files, default: Asci)
\usepackage{lmodern} % vectorized version of the "computer modern" font (default) (https://tex.stackexchange.com/q/1390/148164, https://tex.stackexchange.com/q/147194/148164)
\usepackage[ngerman, english]{babel} %% last language is main
%\shorthandoff{"} %babel breaks tikz-cd otherwise (activate if used)
\usepackage{csquotes} %% biblatex wants csquotes (localized quotes) if babel is present

%%% GENERAL IMPROVEMENTS
\usepackage[unicode, colorlinks]{hyperref} %% clickable links
\usepackage{enumitem} %% better (customizable) enumerations (e.g. \begin{enumerate}[label=(\roman*)])


%%% STYLING
% \setlength{\parindent}{0em} % No indentation for new paragraphs
\usepackage{emptypage} % removes page numbers on empty pages
% \makeatletter \@twosidefalse \makeatother
\usepackage{xcolor}
\usepackage[many]{tcolorbox}
\usepackage{fancyhdr}
\usepackage{mdframed}
\usepackage{booktabs}
\usepackage{multirow}

\usepackage{marvosym} % special letters


%%% GRAPHICS & PLOTS

% \usepackage{graphicx} % enable pictures
\usepackage{wrapfig} % text wrapping around figures
\usepackage{wrapstuff}
\usepackage{tikz, tikz-cd, pgfplots} % plots and sketches
\usetikzlibrary{bending,positioning,calc}
\usetikzlibrary{decorations.pathreplacing}
\usetikzlibrary{arrows.meta}
\pgfplotsset{compat=1.18}


\usepackage{subcaption}

%%% CODE
\usepackage[ruled,linesnumbered]{algorithm2e}

% fix algomathdisplay (see: https://tex.stackexchange.com/a/611426/148164)
\makeatletter
\renewenvironment{algomathdisplay}
 {\[}
 {\@endalgocfline\vspace{-\baselineskip}\]{\DontPrintSemicolon\;}}
\makeatother
%%% MATH
\usepackage{mathtools} % fixes some things from amsmath and adds some new features
\usepackage{amsthm, amssymb} % amsthm -> proof env, amssymb -> additional symbols
% \usepackage{thmtools} % restateable theorem environments
\usepackage{aligned-overset} % https://tex.stackexchange.com/questions/257529/overset-and-align-environment-how-to-get-correct-alignment
%% usage: \overset{something}&{=}
\usepackage{xfrac}

%%% Special Symbols
% \usepackage{bbm} % blackboard numbers: usage \mathbbm{1} (amssymb \mathbb{1} looks bad)
% \usepackage{cancel} % strike through equation parts to indicate they cancel out
% %% contradiction symbol
% \usepackage{wasysym} % lightning symbol
% \newcommand*{\contradiction}[1][]{\quad\text{{\LARGE\lightning} #1}\ }

%%% BIBLIOGRAPHY
\usepackage[numbers]{natbib}

% Editing
\usepackage[draft]{fixme}
\fxsetup{theme=color}

\usepackage{xsim}


%% AMS Theorem Environments
\theoremstyle{plain}% default
\newtheorem{proposition}{Proposition}[section]
\newtheorem{lemma}[proposition]{Lemma}
\newtheorem{corollary}[proposition]{Corollary}
\newtheorem{theorem}[proposition]{Theorem}
\newtheorem{fact}[proposition]{Fact}

\theoremstyle{definition}
\newtheorem{definition}[proposition]{Definition}
\newtheorem{example}[proposition]{Example}
\newtheorem{assumption}[proposition]{Assumption}

\theoremstyle{remark}
\newtheorem{remark}[proposition]{Remark}
%%

\newcommand*{\dims}{d}

\DeclarePairedDelimiter{\abs}{\lvert}{\rvert}
\DeclarePairedDelimiter{\Set}{\{}{\}}
\NewDocumentCommand{\Indicator}{s m}{%
  \IfBooleanTF{#1}
    {\mathbf{1}_{\left\{#2\right\}}} % starred: wrap in braces
    {\mathbf{1}_{#2}}                 % unstarred: simple subscript
}

% spaces
\newcommand*{\real}{\mathbb{R}}
\newcommand*{\nat}{\mathbb{N}}
\newcommand*{\integer}{\mathbb{Z}}
\newcommand*{\complex}{\mathbb{C}}
\newcommand*{\rational}{\mathbb{Q}}

% Topology
\newcommand*{\closure}[1]{\overline{#1}}

% Probability
\newcommand{\E}{\mathbb{E}}
\renewcommand{\Pr}{\mathbb{P}}
\DeclareMathOperator{\Var}{Var}
\DeclareMathOperator{\Cov}{Cov}
% \newcommand*{\normal}{\mathcal{N}}
\newcommand{\normal}[1]{{\color{blue}\boldsymbol{\mathcal{N}}}\left(#1\right)}
\newcommand*{\density}{\varphi}
\DeclareMathOperator{\iid}{iid}
\newcommand*{\given}{\mid}

\newcommand*{\Binomial}{\mathrm{Bin}}
\newcommand*{\Poisson}{\mathrm{Poi}}

% random functions
\newcommand*{\rf}{\mathbf{f}}

% Linear Algebra
\DeclareMathOperator{\sgn}{sgn}
\DeclareMathOperator{\Span}{span}
\newcommand{\identity}{\mathbb{I}}
\newcommand*{\transpose}{T}

% Optimization
\DeclareMathOperator*{\argmin}{arg\,min}
\DeclareMathOperator*{\argmax}{arg\,max}

% Machine Learning
\newcommand*{\cost}{J}
\newcommand*{\Cost}{\mathbf{J}}
\newcommand*{\loss}{\ell}
\newcommand*{\noise}{\varsigma}
\newcommand*{\param}{w}
\newcommand*{\model}{\phi}
\newcommand*{\X}{X}
\newcommand*{\Y}{Y}


% Numerical Analysis
\newcommand*{\bigO}{\mathcal{O}}


%% =========== Color (use RGB for digital, CMYK for print) ==========

% ----- Uni Mannheim Corporate Design -----
\definecolor{primary}{RGB}{0,48,86}
% \definecolor{primary}{cmyk}{1,0.6,0.1,0.6}

\colorlet{primary10}{primary!10}
\colorlet{primary25}{primary!25}
\colorlet{primary40}{primary!40}
\colorlet{primary55}{primary!55}
\colorlet{primary70}{primary!70}
\colorlet{primary85}{primary!85}

\definecolor{accent}{RGB}{179,182,185}
% \definecolor{accent}{cmyk}{35,25,25,0}

% ------------------------------------------

% --- convenient color access ---
\newcommand*{\colorPrimary}[2][]{{\color{primary#1} #2}}
\newcommand*{\colorAccent}[1]{{\color{accent} #1}}

 % AFTER the solution true/false macro !!

% \editmodetrue

\ifeditmode
  \NewDocumentEnvironment{solution}{O{Solution} +b}{\begin{proof}[#1] #2\end{proof}}{}
  \newtheorem{exercise}{Exercise}[section]
\else
  \usepackage{xsim}
  \DeclareExerciseEnvironmentTemplate{theorem}{%
    \par\addvspace{10pt}
    \noindent\textbf{%
      \XSIMmixedcase{\GetExerciseName}~\GetExerciseProperty{counter}%
    }{\normalfont%
      \GetExercisePropertyT{subtitle}{ (\PropertyValue)}}\textbf{. }
    \IfInsideSolutionF{%
      \GetExercisePropertyT{points}{%
        \marginpar{%
          \printgoal{\PropertyValue}%
          \GetExercisePropertyT{bonus-points}{+\printgoal{\PropertyValue}}%
          \,\IfExerciseGoalSingularTF{points}
            {\XSIMtranslate{point}}
            {\XSIMtranslate{points}}%
        }%
      }%
    }%
  }{\par\addvspace{\baselineskip}}

  \renewcommand\printsolutions{%
    \def\currentchapter{}%
    \def\currentsection{}%
    \def\lastchapter{}%
    \def\lastsection{}%
    \chapter*{Solutions\addcontentsline{toc}{chapter}{Solutions}}%
    
    \ForEachUsedExerciseByType{%
      \let\lastchapter\currentchapter
      \let\lastsection\currentsection
      \edef\currentchapter{\ExercisePropertyGet{##1}{##2}{chapter-value}}%
      \edef\currentsection{\ExercisePropertyGet{##1}{##2}{section-value}}%
      % \ifx\lastchapter\currentchapter\else
      %   \addcontentsline{toc}{chapter}{Solutions Chapter \ExercisePropertyGet{##1}{##2}{chapter}}
      %   \chapter*{Solutions Chapter \ExercisePropertyGet{##1}{##2}{chapter}}
      % \fi
      \ifx\lastsection\currentsection\else
        \section*{Exercises \ExercisePropertyGet{##1}{##2}{section}}%
        \addcontentsline{toc}{section}{Exercises \ExercisePropertyGet{##1}{##2}{section}}%
      \fi
      \XSIMprint{solution}{##1}{##2}%
    }%
  }
  \xsimsetup{
      exercise/template=theorem,
      solution/template=theorem,
      exercise/the-counter = \thesection.\arabic{exercise},
      exercise/within=section,
      % solution/print=true,
  }
  % \SetExerciseParameters{exercise}{counter=theorem}
\fi
  \xsimsetup{
      exercise/the-counter = \arabic{exercise},
      solution/print=true,
  }


% Meta Information %%%%%%%%%%%%%%%%%%%%%%%%%%%%%%%%%%%

% included in all sheets (-> only one update per semester)
\course{Compl\'ements de probabilit\'es et statistiques}
\semester{Spring 2026}
\author{Felix Benning}

 % (file) for one update per semester
\sheetNumber{3}

%%%% Document %%%%%%%%%%%%%%%%%%%%%%%%%%%%%%%%%%%%%%%%
\begin{document}

\maketitle

\textbf{
The goal is to have something to talk about in the oral exam.  Skip an exercise
if you get stuck. That is why there are many.
% I am more interested in how you
% would solve them than the actual solution, i.e.\ do not prioritize tedious
% calculations.
}

\begin{exercise}[subtitle={Convergence in distribution but not probability}]
    \label{ex: convergence in distribution but not in probability}
    Let \(X\) be a random variable with \(\prob{X=-1} = \frac12\) and \(\prob{X=1}=\frac12\).
    \[
        X_n \coloneq (-1)^n X
    \]
    Prove that
    \begin{enumerate}[label=(\alph*), noitemsep]
        \item\label{it: (conv in dist but not in prob) in dist} \(X_n \xrightarrow{d} X\)
        \item \label{it: (conv in dist but not in prob) not in prob}\(X_n\) does not converge in probability to \(X\).
    \end{enumerate}
\end{exercise}

\begin{solution}
\begin{enumerate}[wide=0pt]
    \item[\ref{it: (conv in dist but not in prob) in dist}]
    Observe that the distribution of \(X\) is symmetric, that is
    \[
        \prob{-X \le x} = \prob{X \le x} \qquad \forall x\in \real.
    \]
    Consequently, for any \(x \in \real\) we have \(\prob{X_n \le x} = \prob{X \le x}\).
    In particular this implies \(\prob{X_n \le x} \to \prob{X \le x}\) for all \(x\).

    \item[\ref{it: (conv in dist but not in prob) not in prob}]
    We have for \(\epsilon = 1\) and every odd \(n\)
    \[
        \prob{\abs{X - X_n} \ge \epsilon}
        = \prob{2\abs{X} \ge 1} = 1
    \]
    The sequence \(\prob{|X_n - X| \ge \epsilon}\) thus does not converge to \(0\).
    \qedhere
\end{enumerate}
\end{solution}




\begin{exercise}[subtitle={Borel's \(0\)-\(1\) law}]
    \label{ex: borel cantelli 0-1 law}
    Let \((A_n)_{n=1}^\infty\) be a sequence of independent events. Prove that 
    the probability of their liminf and limsup are either zero or one, that is
    \[
        \prob[\Big]{\limsup_{n\to \infty} A_n} \in \{0,1\}
        \quad\text{and}\quad
        \prob[\Big]{\liminf_{n\to \infty} A_n} \in \{0,1\}.
    \]
    Consequently, if the probability is greater zero it must be one!
\end{exercise}

\begin{solution}
    Observe that there are only two possibilities, either
    \[
        \sum_{n=1}^\infty \prob{A_n} < \infty
            \quad\text{or}\quad
        \sum_{n=1}^\infty \prob{A_n} = \infty.
    \]
    In the first case we have by the first Borel-Cantelli lemma \eqref{eq: borel cantelli 1} that
    \[
        \prob[\Big]{\limsup_{n\to \infty} A_n} = 0.
    \]
    in the second case we have by the second Borel-Cantelli lemma \eqref{eq: borel cantelli 2} that
    \[
        \prob[\Big]{\limsup_{n\to \infty} A_n} = 1.
    \]
    Consequently, we have the claim for \(\limsup_n A_n\). Due to
    \[
        \liminf_n A_n = \limsup_n A_n^\complement
    \]
    and the independence of
    \(A_n^\complement\), the
    same arguments apply to \(A_n^\complement\) so we also have the claim for
    \(\liminf_n A_n\).
    \qedhere
\end{solution}



\begin{exercise}[subtitle={Delta method}]
    \label{ex: delta method}
    Let \(X_n\) be a sequence of random variables in \(\real\) such that
    \[
        \sqrt{n}(X_n - \mu) \overset{d}\to \normal{0, \sigma^2}
    \]
    for some \(\mu\in \real\) and \(\sigma^2 > 0\). Let \(g\colon\real\to\real\)
    be a function that is continuously differentiable at \(\mu\) with \(g'(\mu)\neq 0\).
    We will prove the \textbf{delta method} together, which states
    \begin{equation}
        \label{eq: delta method}    
        \sqrt{n}(g(X_n) - g(\mu)) \overset{d}\to \normal[\big]{0, (g'(\mu))^2\sigma^2}.
    \end{equation}
    The following steps will guide you through the proof:
    \begin{enumerate}[label=(\alph*), noitemsep]
        \item\label{it: (delta-method) X_n to mu} Prove that \(X_n \overset{p}\to \mu\).
        \item\label{it: (delta-method) taylor expansion} For the mean value Taylor expansion at \(\mu\)
        \[
            g(X_n) = g(\mu) + g'(\xi_n)(X_n - \mu)
        \]
        with \(\xi_n \in \set{\lambda X_n + (1-\lambda)\mu : \lambda \in [0,1]}\)
        prove that \(\xi_n \overset{p}\to \mu\) and deduce that \(g'(\xi_n) \overset{p}\to g'(\mu)\).
        \item\label{it: (delta-method) proof} Prove the delta method \eqref{eq: delta method}.
    \end{enumerate}
\end{exercise}
\begin{solution}
    \begin{enumerate}[wide=0pt]
        \item[\ref{it: (delta-method) X_n to mu}] For \(N\sim \normal{0,\sigma^2}\) we have 
        by Slutzky
        and continuous mapping that
        \[
            X_n -\mu
            = \underbrace{\tfrac1{\sqrt{n}}}_{\overset{p}\to 0}
            \underbrace{\sqrt{n}(X_n - \mu)}_{\overset{d}\to N}
            \overset{d}\to 0 \cdot N = 0,
        \]
        Since convergence in distribution to a constant implies convergence in probability
        we have \(X_n - \mu\overset{p}\to 0\) and consequently \(\X_n \overset{p}\to \mu\).

        \item[\ref{it: (delta-method) taylor expansion}] For any \(n\in \nat\) we
        have
        \[
            \xi_n = \lambda_n X_n + (1-\lambda_n)\mu
        \]
        for a (random) \(\lambda_n \in [0,1]\). Since \(X_n \overset{p}\to
        \mu\) this implies 
        \[
            \norm{\xi_n - \mu}
            = \abs{\lambda_n} \norm{X_n - \mu}
            \le \norm{X_n - \mu}
            \overset{p}\to 0.
        \]
        Since \(g'\) is continuous at \(\mu\) we have the claim \(g'(\xi_n)
        \overset{p}\to g'(\mu)\) by the continuous mapping.
        
        \item[\ref{it: (delta-method) proof}] Using the mean value Taylor
        expansion we have for \(N\sim \normal{0, \sigma^2}\)
        by Slutzky and continuous mapping that
        \[
            \sqrt{n}(g(X_n) - g(\mu))
            = \underbrace{g'(\xi_n)}_{\overset{p}\to g'(\mu)} \underbrace{\sqrt{n}(X_n - \mu)}_{\overset{d}\to N}
            \overset{d}\to g'(\mu) N \sim \normal[\big]{0, (g'(\mu))^2\sigma^2}. 
            \qedhere
        \]
    \end{enumerate}
\end{solution}


\begin{exercise}[subtitle={Lévy's continuity theorem in infinite dimensions}]
    \label{ex: counterexample levy's continuity in infinite dim}
    To show that Lévy's continuity theorem does not hold in infinite dimensions
    consider the deterministic sequence \(X_n = e^n\), where \(e^n=(0,\dots,0, 1, 0, \dots)\) is the
    \(n\)-th unit vector in the Hilbertspace of square summable sequences 
    \[
        \ell^2 = \set[\Big]{(x_n)_{n\in \nat} \mid \sum_{n=1}^\infty x_n^2 < \infty}
        \qquad\text{with}\qquad
        \scp{x, y} = \sum_{n=1}^\infty x_n y_n.
    \]
    Let \(X=0\in \ell^2\) be the deterministic zero vector. Show that
    \begin{enumerate}[label=(\alph*), noitemsep]
        \item\label{it: (counterexample levy's continuity in infinite dim) charfct}
        \(\charFct_{X_n}(h) \to \charFct_{X}(h)\) for all \(h\in \ell^2\),
        \item\label{it: (counterexample levy's continuity in infinite dim) convergence in distribution}
        \(X_n \not\overset{d}\to X\).
    \end{enumerate}
\end{exercise}

\begin{solution}
\begin{enumerate}[wide=0pt, noitemsep]
    \item[\ref{it: (counterexample levy's continuity in infinite dim) charfct}]
    For any \(h\in \ell^2\) we have \(h_n \to 0\) and consequently
    \[
        \scp{h, e^n} = h_n \to 0.
    \]
    By definition of the characteristic function we thus have
    \[
        \charFct_{X_n}(h)
        = \E{e^{i\scp{h, X_n}}}
        = e^{i\scp{h, e^n}}
        = e^{i h_n}
        \overset{n\to\infty}\to 1
        = e^{i\scp{h, 0}}
        = \charFct_X(h).
    \]
    \item[\ref{it: (counterexample levy's continuity in infinite dim) convergence in distribution}]
    Since \(X\) is deterministic, \(X_n \overset{d}\to X\) would imply \(X_n
    \overset{p}\to X\).
    And since \(\norm{X_n} = 1\) for all \(n\) the \(X_n, X\) are bounded
    which would imply
    convergence in \(L^2\). But
    \[
        \E{\norm{X_n - X}^2}
        = \norm{e^n}^2
        = 1
        \not\to 0.
        \qedhere
    \]
\end{enumerate}    
\end{solution}


% \begin{exercise}[subtitle={Infinitely often or never successful}]
%     Let \(X_n \sim \bernoulli(n^{-r})\) be independent Bernoulli random variables,
%     whose success probability decreases in \(n\) at rate \(r > 0\). Note that we say
%     that \(X_n\) is a success if \(X_n = 1\). Show that
%     \begin{enumerate}[label=(\alph*)]
%         \item\label{it: finitely many successes}
%         for \(r > 1\) there are finitely many successes with probability one.
%         That is after some random but finite time \(N\in \nat\) all \(X_n\) are zero.

%         \item\label{it: infinitely many successes}
%         for \(r \le 1\) there are infinitely many successes with probability one.
%     \end{enumerate}
% \end{exercise}

% \begin{solution}
% \begin{enumerate}[wide=0pt]
%     \item[\ref{it: finitely many successes}]
%     Since we have
%     \[
%         \sum_{n=1}^\infty \prob{X_n = 1} = \sum_{n=1}^\infty n^{-r} < \infty,
%     \]
%     we have by the first Borel-Cantelli lemma \eqref{eq: borel cantelli 1} that
%     \[
%         \prob{X_n=1 \text{ for infinitely many } n} = \prob{\limsup_{n\to \infty} \set{X_n = 1}} = 0.
%     \]

%     \item[\ref{it: infinitely many successes}]
%     Since we have
%     \[
%         \sum_{n=1}^\infty \prob{X_n = 1} = \sum_{n=1}^\infty n^{-r} = \infty,
%     \]
%     we have by the second Borel-Cantelli lemma \eqref{eq: borel cantelli 2} that
%     \[
%         \prob{X_n=1 \text{ for infinitely many } n} = \prob{\limsup_{n\to \infty} \set{X_n = 1}} = 1.
%         \qedhere
%     \]
% \end{enumerate}
% \end{solution}

\begin{exercise}[subtitle={Single issue polling}]
    \label{ex: single issue polling} 
    Let \(x_1,\ldots, x_N \in \set{\text{no}, \text{yes}} \simeq \set{0,1}\) be the
    opinions of \(N\) people on a certain decision. In the following we will
    assume \(x_i \in \set{0,1}\). Let
    \[
        p = \frac1N\sum_{i=1}^N x_i
    \]
    be the share of people that agree with the decision.

    Assume that you can randomly sample \(n\) people from the population and
    and that they truthfully tell you their opinions \(X_1, \dots, X_n\).  That
    is \(\prob{X_i = x_j} = \frac1N\) for \(j=1, \ldots, N\) and the \(X_i\) are
    independent.
    \begin{enumerate}[label=(\alph*), noitemsep]
        \item\label{it: (single issue polling) distribution}
         Prove the \(X_i\) are Bernoulli distributed, specifically \(X_i \sim \bernoulli(p)\). 

        \item\label{it: (single issue polling) PAC bound}
        Prove that
        \[
            \prob[\Big]{\abs[\Big]{\frac1n \sum_{i=1}^n X_i - p} > \epsilon}
            \le 2\exp(-2n \epsilon^2).
        \]
    \end{enumerate}
\end{exercise}

\begin{solution}
\begin{enumerate}[wide=0pt]
    \item[\ref{it: (single issue polling) distribution}]
    We have
    \begin{align*}
        \prob{X_i = 1}
        &= \sum_{j=1}^N \prob{X_i = 1, X_i = x_j}
        \\
        &= \sum_{j=1}^N \prob{X_i = 1 | X_i = x_j} \prob{X_i = x_j}
        = \sum_{j=1}^N x_j \frac1N = p.
    \end{align*}

    \item[\ref{it: (single issue polling) PAC bound}]
    Since \(X_i \in \set{0,1}\) we can apply Hoeffding's inequality for
    bounded random variables to get the claim.

    \qedhere
\end{enumerate} 
\end{solution}


\end{document} 

%%% Local Variables: 
%%% mode: latex
%%% TeX-master: t
%%% End: 

